\section{Multimodal fusion and executable grounding}
\label{sec:C05:fusion}
A multimodal language model commonly maps visual features through a projection or adapter into representations consumed by a language decoder. Visual instruction tuning studies how such a system can answer or follow instructions involving images~\cite{liu2023llava}. For an agent, the output must eventually identify an executable target and arguments. This adds a grounding obligation absent from many question-answering tasks: the referenced object must correspond to the current environment, and the selected action must be valid there. Describing the correct button is not equivalent to clicking it, and producing a plausible screenshot description is not evidence that a state transition occurred.

Fusion strategies expose different bottlenecks. Early token fusion lets attention combine modalities directly but can consume a large context budget. A fixed-size visual summary controls cost while restricting recoverable detail. Cross-attention can separate language and visual lengths but adds its own compute and architectural assumptions. A tool-mediated approach leaves images outside the main model and retrieves structured descriptions or crops on demand. None dominates independently of task distribution. Exact document reading, spatial manipulation, and long video understanding demand different balances of detail, temporal coverage, and interaction latency.

The data pipeline must preserve the relation between natural-language instructions and observations. If a training example includes a final answer beside a screenshot taken after the action, the learner may exploit evidence unavailable when the action was chosen. A similar leak occurs when an evaluator supplies a full trajectory while judging a decision that should depend only on its prefix. Dataset construction should retain action timestamps and distinguish pre-action observations from post-action verification. The context-policy consistency issue is thus not restricted to text compression; it also applies to multimodal temporal alignment.

Untrusted visual content can contain instructions. A web page, document, or terminal screenshot can ask the agent to reveal secrets or change its objective. Converting that content to text through OCR does not change its authority. The system should retain provenance through every modality transformation and enforce authorization outside the model. A learned refusal policy can reduce harmful routing, but cannot replace an execution boundary that checks actual permissions. Evaluations should include task-relevant visual instructions and adversarial embedded instructions so that ignoring all image text is not rewarded as a successful defense.
